%%%%%%%%%%%%%%%%%%%%%%%%%%%%%%%%%%%%%%%%%%%%%%%%%%%%%%%%%%%%%%%%%%%%%%%%%%%%%%%%
%2345678901234567890123456789012345678901234567890123456789012345678901234567890
%        1         2         3         4         5         6         7         8

\documentclass[a4paper, 10pt, conference]{ieeeconf}

\IEEEoverridecommandlockouts

\overrideIEEEmargins

\usepackage[utf8]{inputenc}
\usepackage[T1]{fontenc}
\usepackage[vietnamese]{babel}
\usepackage{graphicx}
\usepackage{amsmath}

%%%%%%%%%%%%% added for project course control engineering
\usepackage[left=1.57cm,right=1.57cm,top=2.54cm,bottom=2.54cm]{geometry}
\usepackage{subfigure}
%%%%%%%%%%%%% added for project course control engineering


\usepackage{amssymb}
\usepackage{enumitem}
\usepackage{multirow}
\usepackage{amsthm}

%%%%%%%%%%%%% added for the Thermal Digital Twin (DT4TM) report %%%%%%%%%%%%%
\usepackage{booktabs}
\usepackage{tikz}
\usetikzlibrary{arrows.meta, positioning}
\usepackage{url}
\usepackage{tabularx}
\usepackage[hidelinks]{hyperref}
%%%%%%%%%%%%% added for the Thermal Digital Twin (DT4TM) report %%%%%%%%%%%%%

\theoremstyle{remark}
\newtheorem{remark}{Ghi chú}

\makeatletter
\def\@IEEEprocessthesectionargument#1{%
\@ifmtarg{#1}{%
\@IEEEappendixsavesection*{Phụ lục \thesectiondis}%
\addcontentsline{toc}{section}{Phụ lục \thesection}}{%
\@IEEEappendixsavesection*{Phụ lục \thesectiondis \\* #1}%
\addcontentsline{toc}{section}{Phụ lục \thesection: #1}}}
\makeatother

\title{\LARGE \bf
Cặp Kỹ Thuật Số Nhiệt cho Hệ Thống Lơ Lửng Điện Động
}
\author{ Dang Minh Hoang, Marina Borchert, Mohamed Aziz El Majid
\thanks{Khóa học Dự án Kỹ Thuật Số, Học kỳ Hè 2026. Người hướng dẫn: GS. TS. Dirk Hartmann, PGS. TS. Melina Merkel}
}


\begin{document}



\maketitle
\thispagestyle{plain}
\pagestyle{plain}


%%%%%%%%%%%%%%%%%%%%%%%%%%%%%%%%%%%%%%%%%%%%%%%%%%%%%%%%%%%%%%%%%%%%%%%%%%%%%%%%
\begin{abstract}
Báo cáo này mô tả khái niệm, phương pháp luận và kiến trúc phần mềm của một
\emph{kỹ thuật số} (Digital Twin) để dự đoán nhiệt trên máy lơ lửng điện động
theo chuẩn TEAM-28 (COMPUMAG). Thay vì sử dụng phần mềm CAE thương mại, chúng tôi
chọn một trình giải Python tự phát triển, vì mục tiêu thực sự không phải là một
lần mô phỏng mà là một \emph{mô hình dự đoán thời gian thực} chạy trên Điện thoại
thông minh thông qua chế độ Thực tế Tăng cường (Augmented Reality). Báo cáo này
giải thích lựa chọn này từ quan điểm vật lý (tổn thất Joule tỷ lệ với bình phương
dòng điện với phân bố không gian cố định, cho phép một lần giải FEM ngoại tuyến
cộng với một mô hình giảm bậc thời gian thực). Dòng điện cuộn dây \emph{được đo}
$I(t)$ đóng vai trò là đầu vào: cảm biến Hall trên giá thử nạp liên tục vào bộ
tích phân thời gian, trong khi nhiệt độ là đầu ra được tính toán từ đó. Báo cáo
tóm tắt các kết quả trước đây và dự kiến, đồng thời phân biệt rõ ràng giữa trạng
thái ổn định đã được xác thực và quá trình thời gian chưa được xác thực. Để triển
khai mô hình vật lý, chúng tôi sử dụng một công cụ AI được tích hợp trong môi
trường phát triển, trong khi tất cả các giả định vật lý và tiêu chí xác thực được
định nghĩa và kiểm tra bởi nhóm.
\end{abstract}

%%%%%%%%%%%%%%%%%%%%%%%%%%%%%%%%%%%%%%%%%%%%%%%%%%%%%%%%%%%%%%%%%%%%%%%%%%%%%%%


\section{Giới thiệu}
\label{sec:intro}

\subsection{Mô tả công việc}
Vấn đề TEAM~28~\cite{team28} là một chuẩn mực công nhân được thiết lập của cộng
đồng Điện từ Tính toán (COMPUMAG) để xác thực các trình giải điện từ tự viết: Một
đĩa nhôm dẫn điện lơ lửng phía trên một bộ cuộn dây được chạy bằng dòng xoay chiều
($I_{\mathrm{rms}}\approx5\,\mathrm{A}$ được đo, $f=50\,\mathrm{Hz}$). Từ thông
thay đổi theo thời gian tạo ra các dòng xoáy trong đĩa, tạo ra cả lực nâng (lơ
lửng) và mất mát nhiệt. Giá thử nghiệm được xây dựng tại khoa có cấu trúc tương
tự với chuẩn mực này, nhưng có thêm một bộ cuộn dây được chia thành hai phần
(cuộn trong 1000 vòng, cuộn ngoài 500 vòng) với lõi từ ferromạgnetic
($\mu_r\approx1000$) và vành ngoài ferromạgnetic, không xuất hiện trong vấn đề ban
đầu và tự tạo ra nhiệt đáng kể.

\begin{figure}[t]
\centering
\begin{tikzpicture}[x=0.065cm, y=0.065cm,
    ironfill/.style={fill=gray!45},
    cufill/.style={fill=orange!35},
    alfill/.style={fill=cyan!25},
    lbl/.style={font=\scriptsize, align=center}]
  % Symmetry axis
  \draw[dash dot] (0,15) -- (0,-58);
  \node[font=\tiny, right] at (2,16) {$r=0$};
  % r-axis
  \draw[-{Stealth[length=4pt]}] (0,-58) -- (109,-58) node[below, font=\tiny] {$r$ [mm]};
  \draw (25.9,-58) -- (25.9,-56.5); \node[font=\tiny] at (25.9,-63) {25{,}9};
  \draw (61.9,-58) -- (61.9,-56.5); \node[font=\tiny] at (61.9,-63) {61{,}9};
  \draw (79.9,-58) -- (79.9,-56.5); \node[font=\tiny] at (79.9,-63) {79{,}9};
  \draw (102.9,-58) -- (102.9,-56.5); \node[font=\tiny] at (102.9,-63) {102{,}9};
  % Iron core
  \draw[ironfill] (0,0) rectangle (25.9,-53);
  \node[lbl] at (13,-26.5) {Lõi sắt\\$\mu_r{=}1000$};
  % Inner coil
  \draw[cufill] (27.9,0) rectangle (61.9,-53);
  \node[lbl] at (44.9,-26.5) {Cuộn trong\\$N{=}1000$};
  % Iron ring
  \draw[ironfill] (64.9,0) rectangle (79.9,-53);
  \node[lbl, rotate=90] at (72.4,-26.5) {Vành sắt};
  % Outer coil
  \draw[cufill] (82.9,0) rectangle (102.9,-53);
  \node[lbl] at (92.9,-26.5) {Cuộn\\ngoài\\$N{=}500$};
  % Al-Disc (schematic position)
  \draw[alfill] (0,6) rectangle (80,9);
  \node[lbl] at (40,7.5) {Đĩa Al, $R{=}80$, $3$\,mm};
  \draw[<->, dashed] (85,0) -- (85,6);
  \node[font=\tiny] at (91,3) {$z_{\mathrm{gap}}$};
\end{tikzpicture}
\caption{Mặt cắt ngang của hình học thiết bị, giống với đầu vào trong
\texttt{params.yaml}. Đầu vào $I_{\mathrm{rms}}=5\,\mathrm A$: được tính toán
$z_{\mathrm{gap}}\approx12\,\mathrm{mm}$, thực tế quan sát được $7$--$8\,\mathrm{mm}$
(Phần~\ref{sec:limitations}).}
\label{fig:geometry}
\end{figure}

\begin{table}[t]
\centering\footnotesize
\caption{Hình học thiết bị thực tế được đo bằng thước (2026-07-10) -- dữ liệu đầu
vào chính xác cho trình giải FEM (\texttt{params.yaml}), có thể được sử dụng để
so sánh trực tiếp với chuẩn gốc hoặc để xây dựng lại trong phần mềm CAE thương mại
(Phần~\ref{sec:tools}).}
\label{tab:geometry}
\begin{tabular}{@{}lrrrr@{}}
\toprule
Bộ phận & $r$\,[mm] & $\mu_r$ & $\sigma$\,[MS/m] & $N$ \\
\midrule
Lõi sắt & 0{,}0--25{,}9 & 1000 & 1{,}0 & -- \\
Khoảng không khí & 25{,}9--27{,}9 & 1 & 0 & -- \\
Cuộn trong (Cu) & 27{,}9--61{,}9 & 1 & 59{,}6 & 1000 \\
Khoảng không khí & 61{,}9--64{,}9 & 1 & 0 & -- \\
Vành sắt & 64{,}9--79{,}9 & 1000 & 1{,}0 & -- \\
Khoảng không khí & 79{,}9--82{,}9 & 1 & 0 & -- \\
Cuộn ngoài (Cu) & 82{,}9--102{,}9 & 1 & 59{,}6 & 500 \\
Đĩa Al & 0{,}0--80{,}0 & 1 & 34{,}0 & -- \\
Không khí $\to$ khung & 102{,}9--130{,}4 & 1 & 0 & -- \\
Tường khung (gỗ ép) & 130{,}4--152{,}9 & 1 & 0 & -- \\
\bottomrule
\end{tabular}

\vspace{2pt}
{\footnotesize Chiều cao của tất cả các vùng lõi/cuộn/vành: $53\,\mathrm{mm}$;
độ dày đĩa: $3\,\mathrm{mm}$ (được định vị phía trên chiều cao này, so sánh
Hình~\ref{fig:geometry}). Kích thích điện: $I_{\mathrm{rms}}=5{,}0\,\mathrm A$
(điểm hoạt động chính, phím quay Variac $220^\circ\approx195{,}5\,\mathrm V$)
hoặc $7{,}8\,\mathrm A$ (phím Variac ở vị trí tối đa $270^\circ$, $\approx240\,\mathrm V$,
điểm lập chỉ số hiệu chuẩn), $f=50\,\mathrm{Hz}$.}
\end{table}

Hình~\ref{fig:geometry} và Bảng~\ref{tab:geometry} ghi lại hình học thực tế này
hoàn toàn và giống hệt nhau với \texttt{params.yaml} -- \emph{nguồn duy nhất}
mà từ đó cả trình giải FEM của chính chúng tôi và chạy so sánh tiềm năng trong
phần mềm thương mại được cung cấp. Vì lý do minh bạch: chỉ các bán kính, số vòng
cuộn dây và lớp vật liệu (sắt/đồng/nhôm được xác nhận bằng bài kiểm tra từ tính
ngày 2026-07-10) được đo trực tiếp; độ thẩm thấu $\mu_r=1000$ là một giá trị
trình giữ chỗ không được đo (Phần~\ref{sec:limitations}), và độ dẫn $\sigma$ là
các giá trị từ tài liệu tham khảo cho từng lớp vật liệu, không phải đo trên mẫu
cụ thể.

\subsection{Từ mô hình mô phỏng đến kỹ thuật số}
Nhiệm vụ không phải là tính toán \emph{một} trường nhiệt độ, mà xây dựng một mô hình
(i) dự đoán trường nhiệt độ $T(r,z,t)$ \textbf{liên tục}, được điều khiển bởi dòng
cuộn dây \textbf{được đo} thực sự $I(t)$, (ii) \textbf{xác thực} với giá thử nghiệm
thực tế (Máy ảnh IR), và (iii) cuối cùng có thể được xem qua \textbf{Mã QR} trên
điện thoại thông minh trong Thực tế Tăng cường. Ba yêu cầu này cùng nhau xác định
một \emph{kỹ thuật số} theo đúng nghĩa~\cite{tao2019digitaltwin} và là điểm khởi
đầu của tất cả các quyết định phương pháp luận được giải thích sau.


\section{Khái niệm Vật lý}
\label{sec:concept}

Một kỹ thuật số tính toán lại mô phỏng FEM đầy đủ mỗi khi dòng điện thay đổi sẽ
không thể thực hiện được trên điện thoại thông minh (thời gian tính toán: từ giây
đến phút trên mỗi lần chạy). Giải pháp nằm trong một quan sát vật lý, không phải
một phím tắt phần mềm: Với tần số cố định ($50\,\mathrm{Hz}$) và các tính chất vật
liệu tuyến tính, tất cả các tổn thất Joule tỷ lệ với bình phương dòng điện
($P\propto I^2$), trong khi \textbf{phân bố không gian} của các tổn thất này
\emph{không} thay đổi -- chỉ có biên độ của chúng. Điều này dẫn đến một kiến trúc
hai bước:

\begin{enumerate}[nosep]
  \item \textbf{Ngoại tuyến (một lần):} FEM điện từ được giải một lần ở dòng điện
        tham chiếu $I_{\mathrm{ref}}$ và cung cấp một bản đồ mất mát được phân giải
        không gian $\hat q_{\mathrm{Al}}(r,z)$ cũng như mất mát cuộn dây
        $\hat P_{\mathrm{Cu}}$.
  \item \textbf{Trực tuyến (Millisecond):} Đối với bất kỳ dòng điện $I$ nào và
        nhiệt độ cơ thể $\bar T$ được biết lần cuối cùng, những cái này chỉ được
        nhân với một đại lượng vô hướng.
\end{enumerate}

Điều quan trọng là sự hiệu chỉnh $\sigma(T)$ cho hai cơ thể là \textbf{ngược chiều}
(Phần~\ref{sec:physics}): Tổn thất dòng xoáy của đĩa là $\propto\sigma_{\mathrm{Al}}$,
tổn thất cuộn dây ohmic là $\propto1/\sigma_{\mathrm{Cu}}$. Do đó, các phương trình
thời gian thực là
\begin{align}
  q_{\mathrm{Al}}(r,z;I,\bar T) &= \hat q_{\mathrm{Al}}(r,z)
     \left(\frac{I}{I_{\mathrm{ref}}}\right)^{\!2}
     \frac{\sigma_{\mathrm{Al}}(\bar T)}{\sigma_{\mathrm{Al}}(T_{\mathrm{ref}})},
     \label{eq:realtime}\\[2pt]
  P_{\mathrm{Cu}}(I,\bar T) &= \hat P_{\mathrm{Cu}}
     \left(\frac{I}{I_{\mathrm{ref}}}\right)^{\!2}
     \frac{\sigma_{\mathrm{Cu}}(T_{\mathrm{ref}})}{\sigma_{\mathrm{Cu}}(\bar T)}.
     \label{eq:realtime_coil}
\end{align}
Cả hai đều cung cấp cho một mô hình nhiệt giảm bậc (ROM), cũng chỉ hoạt động với
số học vectơ/vô hướng.

\begin{remark}[Trạng thái triển khai]
\label{rem:sigma}
Phương trình~\eqref{eq:realtime} được triển khai hoàn toàn trong mô hình thời gian
chạy. Phương trình~\eqref{eq:realtime_coil} \emph{không phải là}: Các nút cuộn dây
và sắt chỉ được tỷ lệ ở thời gian chạy với $(I/I_{\mathrm{ref}})^2$. Sự phụ thuộc
nhiệt độ của cuộn dây thay vào đó được chứa ngầm trong các giá trị dẫn điện lưu
thông $hA$ được điều chỉnh với dữ liệu đo được (Phần~\ref{sec:validation}). Tại
điểm hiệu chuẩn, điều này là chính xác, đối với các điểm hoạt động xa xôi, nó là
một đơn giản hóa đã biết.
\end{remark}

Sự phân tách này là cốt lõi phương pháp luận của dự án: "Trí thông minh" nằm hoàn
toàn trong trình giải ngoại tuyến một lần; khả năng thời gian thực phát sinh từ cấu
trúc vật lý của vấn đề (Tuyến tính, tần số cố định), không phải từ vật lý đơn giản
hóa.


\section{Mô hình Vật lý}
\label{sec:physics}

Mô hình bao gồm ba vấn đề con được ghép nối, đối xứng trục, được giải trong 2D
$(r,z)$ và quay đến 3D cho hiển thị. Bản sao đầy đủ nằm trong Phụ lục~\ref{app:derivation}.

\textbf{Điện từ (Phasor):} Vì chỉ mất mát nhiệt được tính trung bình trong chu kỳ
là liên quan đến (theo kiểu so với $50\,\mathrm{Hz}$ chậm) nhiệt học, trường không
được mô phỏng theo từng bước thời gian mà được giải dưới dạng một phasor phức
$A_\varphi(r,z)$ -- công thức được sử dụng ở đây thông qua thế năng vectơ từ tính
tuân theo công trình cổ điển của Bíró và Preis~\cite{biro1989vectorpotential},
\begin{equation}
  -\nabla\!\cdot\!\left(\nu\nabla A_\varphi\right)
  + \nu\frac{A_\varphi}{r^{2}} + j\omega\sigma A_\varphi = J_s,
  \label{eq:em}
\end{equation}
với $\nu=1/(\mu_r\mu_0)$, $\omega=2\pi f$ và mật độ dòng điện được tạo $J_s$
trong các cuộn dây (cuộn trong/ngoài với chiều cuộn đối lập). Trong lõi/vành
ferromạgnetic, $\nu$ giảm mạnh và do đó gói lại từ thông.

\textbf{Nhiệt Joule:} Từ mật độ dòng xoáy được tạo ra $J_e=-j\omega\sigma A_\varphi$,
mật độ mất mát được tính trung bình trong chu kỳ tuân theo
\begin{equation}
  q(r,z) = \frac{|J_e|^{2}}{2\sigma}
         = \tfrac12\,\sigma\,\omega^{2}\,|A_\varphi|^{2}\ \;[\mathrm{W/m^3}],
  \label{eq:loss}
\end{equation}
có sự phụ thuộc bình phương của $|A_\varphi|$ là cơ sở cho tỷ lệ $I^2$ từ Phương
trình~\eqref{eq:realtime}.

\textbf{Độ dẫn phụ thuộc nhiệt độ:} Vì $70\,\mathrm K$ gia tăng thay đổi lợi
điện trở cụ thể bởi $\approx27\,\%$, chúng tôi có $\sigma(T)=\sigma_0/(1+\alpha(T-T_0))$
với $\alpha\approx3{,}9\times10^{-3}\,\mathrm{K^{-1}}$ -- với hiệu ứng ngược chiều:
tổn thất dòng xoáy đĩa $\propto\sigma_{\mathrm{Al}}$ (phản hồi âm), tổn thất cuộn
dây $\propto1/\sigma_{\mathrm{Cu}}$ (phản hồi dương) -- xem Phương trình~\eqref{eq:realtime}
và \eqref{eq:realtime_coil} cũng như Ghi chú~\ref{rem:sigma} trên trạng thái triển
khai. Cả hai hiệu ứng chỉ thay đổi biên độ, do đó kiến trúc hai bước vẫn có giá trị.

\textbf{Dẫn nhiệt:} Mật độ mất mát $q$ cung cấp cho phương trình Fourier không
ổn định
\begin{equation}
  \rho c_p\,\frac{\partial T}{\partial t} = \nabla\!\cdot\!(k\nabla T) + q,
  \label{eq:fourier}
\end{equation}
với điều kiện biên Robin $-k\,\partial T/\partial n=h(T-T_\infty)$. Trong trường
hợp ổn định, cân bằng năng lượng $\int_V q\,\mathrm dV=\oint_\Gamma h(T-T_\infty)\,\mathrm dA$
phải được thỏa mãn chính xác -- bài kiểm tra xác minh trung tâm của trình giải
nhiệt (xem Phần~\ref{sec:validation}).


\section{Triển khai Số học và Kiến trúc Phần mềm}
\label{sec:numerics}

Cả hai vấn đề con đều được rời rạc hóa bằng Phương pháp Phần tử Hữu hạn (FEM)
trong công thức Galerkin, với các phần tử tam giác tuyến tính (P1) trên lưới 2D
có cấu trúc $(r,z)$; dạng yếu đối xứng trục mang trọng lượng thể tích $2\pi r$ (chi
tiết trong Phụ lục~\ref{app:derivation}). Lắp ráp điện từ tạo ra một ma trận thưa
phức tạp (vì $j\omega\sigma$), được giải bằng \url{scipy.sparse.linalg.spsolve}
dựa trên NumPy~\cite{harris2020numpy} và SciPy~\cite{virtanen2020scipy}. Với
$\sim\!10^4$ ẩn số (2D, đối xứng trục), một trình giải trực tiếp là đủ; chiều sâu
da trong nhôm ở $50\,\mathrm{Hz}$ ($\delta\approx12\,\mathrm{mm}\gg3\,\mathrm{mm}$
độ dày đĩa) làm cho một lưới tốt ở hướng $z$ là không cần thiết.

Từ lần giải FEM một lần, một mô hình thứ nguyên thấp, khả năng thời gian thực được
rút ra: tỷ lệ $I^2$ với hằng số thời gian bậc nhất ($\tau\approx4{,}07\,\mathrm{min}$
ở $R=80\,\mathrm{mm}$) cho đĩa, cũng như một mạng RC phân tán (\emph{mạng nhiệt
tham số lumped}) cho cuộn dây, lõi/vành sắt và nút không khí chung -- một nguyên
tắc mô hình hóa được thiết lập trong công nghệ điều khiển điện~\cite{wallscheid2021thermal},
có thể được hiệu chỉnh trực tiếp so với dữ liệu cảm biến thực. Toàn bộ bước trực
tuyến giảm thành $O(n)$-số học, khả năng chạy trong mili giây trong trình duyệt của
điện thoại thông minh. Bảng~\ref{tab:modules} (Phụ lục~\ref{app:software}) liệt kê
thêm cho mỗi tệp rõ ràng phương trình nào nó giải và bằng thư viện/phương pháp nào.

\begin{figure*}[t]
\centering
\resizebox{\textwidth}{!}{%
\begin{tikzpicture}[
  node distance=5mm and 9mm,
  box/.style={draw, rounded corners, align=center, minimum width=3.0cm,
              minimum height=0.95cm, font=\footnotesize},
  arr/.style={-Stealth, thick}
]
\node[box] (params) {\texttt{params.yaml}\\[-2pt]\tiny (Nguồn sự thật duy nhất cho tất cả các tham số)};
\node[box, right=of params] (config) {\texttt{config.py}\\[-2pt]\tiny Đơn vị, $I_{\mathrm{peak}}$};
\node[box, right=of config] (em) {\texttt{em\_solver.py}\\[-2pt]\tiny FEM điện từ, Bản đồ mất mát};
\node[box, right=of em] (th) {\texttt{thermal\_solver.py}\\[-2pt]\tiny FEM nhiệt, Cân bằng};
\node[box, right=of th] (rom) {\texttt{rom.py} / \texttt{twin\_core.py}\\[-2pt]\tiny ROM thời gian thực (SSOT)};
\node[box, right=of rom] (out) {\texttt{build\_twin\_html\_fem.py}\\[-2pt]\tiny Nướng $\to$ AR-HTML + QR};
\node[box, below=of rom] (sensor) {\texttt{data\_io.py} + Arduino\\[-2pt]\tiny Dòng điện được đo $I(t)$};

\draw[arr] (params) -- (config);
\draw[arr] (config) -- (em);
\draw[arr] (em) -- (th);
\draw[arr] (th) -- (rom);
\draw[arr] (rom) -- (out);
\draw[arr] (sensor) -- (rom);
\end{tikzpicture}%
}
\caption{Đường ống phần mềm: từ tệp tham số qua các trình giải FEM ngoại tuyến
đến ROM thời gian thực và cung cấp AR cuối cùng. Đường dẫn cảm biến cung cấp dòng
cuộn dây được đo $I(t)$ làm \emph{đầu vào} của bộ tích phân thời gian
(Phần~\ref{sec:realtime_input}); nhiệt độ là đầu ra được tính toán từ nó. Xác
thực chéo (\texttt{xval\_twin.py}) thêm vào đó đảm bảo rằng bộ tích phân thời gian
thực được nướng trong JavaScript phù hợp với $10^{-9}$ (tuyệt đối) với bộ tích phân
tham chiếu Python.}
\label{fig:pipeline}
\end{figure*}

Hình~\ref{fig:pipeline} cho thấy đường ống một chiều: Mỗi tệp có chính xác một
trách nhiệm, và tất cả các tham số vật lý nằm trung tâm trong \texttt{params.yaml}
(\emph{nguồn sự thật duy nhất}) -- không có hằng số nào được dây cứng trong mã.

\subsection{Dòng điện xoay chiều được đo làm Đầu vào Thời gian thực}
\label{sec:realtime_input}

Cho đến nay, kỹ thuật số được vận hành với một giả định \emph{không đổi}
$I=5\,\mathrm A$. Trạng thái mục tiêu thực tế là khác: Dòng cuộn dây nên được
\textbf{đo liên tục} và phản hồi làm đầu vào cho mô hình nhiệt -- chỉ khi đó nó
trở thành kỹ thuật số theo dõi hoạt động thực tế thay vì chỉ mô phỏng nó.

Đo lường, điều này được giải quyết thông qua cảm biến dòng điện tác dụng Hall
(ACS712-20A) trong một dây nối của mạch cuộn dây. Nó được cách ly gal vanically
($\approx2{,}1$\,kV), cần thiết với $\approx195$--$240\,\mathrm V$ điện áp mạng
được áp dụng; một điện trở shunt đơn giản sẽ không được phép ở đây. Bộ điều khiển
vi mô hình thành giá trị hiệu dụng của thành phần xoay chiều xung quanh
\emph{được đo} giá trị trung bình (không phải xung quanh danh nghĩa $V_{\mathrm{cc}}/2$-offset,
vì cả nguồn cấp $5\,\mathrm V$ và điểm không của cảm biến có dung sai và trôi):
\begin{equation}
  I_{\mathrm{rms}} = \frac{k}{S}\sqrt{\frac{1}{N}\sum_{n=1}^{N}
      \bigl(v_n - \bar v\bigr)^2},
  \qquad \bar v = \frac{1}{N}\sum_{n=1}^{N} v_n .
  \label{eq:irms}
\end{equation}
Ở đây $S=0{,}100\,\mathrm{V/A}$ là độ nhạy cảm biến và $k$ là hệ số hiệu chỉnh
một điểm chống lại đồng hồ vôn kế ở điểm hoạt động $5\,\mathrm A$ ($k=1$, miễn là
không được đo); nó hấp thụ dung sai của cả hai hằng số tờ dữ liệu.

Quyết định về độ chính xác là sự lựa chọn của cửa sổ trung bình: Giá trị thực hiệu
dụng chỉ chính xác trong một số \emph{số nguyên} các chu kỳ mạng. Cửa sổ do đó được
xác định trên \textbf{Thời gian} -- năm chu kỳ mạng, đó là chính xác $100\,\mathrm{ms}$ --
và không phải trên số lượng mẫu cố định $N$. Nếu không, kết quả sẽ phụ thuộc vào tốc
độ ADC của bảng cụ thể: cùng $1000$ mẫu kéo dài trên bảng AVR khoảng $5{,}6$ chu kỳ
mạng, trên ADC nhanh hơn của bảng Renesas chỉ khoảng $1{,}25$, gây ra lỗi phụ thuộc
pha từ vài phần trăm. Với cửa sổ giới hạn thời gian, $N$ là một hệ quả của tốc độ
bảng thay vì một nguồn lỗi, và firmware tương tự cung cấp cùng giá trị đo trên cả
hai nền tảng.

Phần mềm, chuỗi được triển khai hoàn toàn: giá trị đo được truyền mỗi giây dưới
dạng CSV và điều khiển bộ tích phân thời gian trực tiếp -- bước thời gian được
thực hiện với $I(t)$ được đo thay vì hằng số (\texttt{data\_io.py} gọi
\texttt{TwinState.step(}$I_{\mathrm{rms}}$\texttt{,}~$\Delta t$\texttt{)} với
$\Delta t$ giữa hai mẫu; khi mẫu không hợp lệ, nó rơi lại vào giá trị hằng số). Vì
tất cả các tổn thất theo Phương trình~\eqref{eq:realtime} tỷ lệ với bình phương
dòng điện, nó vẫn là một phép nhân vô hướng trên mỗi bước thời gian -- khả năng
thời gian thực vẫn được duy trì ngay cả với đầu vào đo lường thực tế. Chuỗi được
kiểm tra hoàn toàn với chuỗi đo lường tổng hợp; chỉ thiếu cài đặt vật lý của cảm
biến.


\section{Tại sao Mô phỏng với Thư viện Python}
\label{sec:tools}

Mô phỏng được tự triển khai trong Python (NumPy~\cite{harris2020numpy},
SciPy~\cite{virtanen2020scipy}) vì yêu cầu thực tế -- một mô hình thời gian thực,
có thể hiệu chỉnh cảm biến, chạy trên điện thoại thông minh -- yêu cầu kiểm soát
ở cấp mã nguồn: Để khai thác cấu trúc $I^2$ từ Phần~\ref{sec:concept} (FEM giải
một lần, sau đó chỉ tỷ lệ), bản đồ mất mát $\hat q(r,z)$ phải được trích xuất
\emph{trước} phép nhân với $I^2$ một cách có chủ đích và tái sử dụng trong một ROM
thời gian thực riêng biệt -- cấu trúc chỉ có thể nếu trình giải tự được viết.

Để định vị, các lựa chọn thay thế được thiết lập được kiểm tra: Các công cụ như
SimScale, COMSOL/ANSYS Maxwell~\cite{comsol_team28} hoặc FEMM~\cite{meeker2020femm}
có thể giải quyết vấn đề điện từ nhiệt một cách cơ bản -- COMSOL thậm chí còn cung
cấp một mô hình ví dụ hoàn thành cho TEAM~28~\cite{comsol_team28} -- cũng như các
trình giải Mã nguồn mở Elmer~\cite{malinen2013elmer} và FEniCS~\cite{langtangen2017fenics}.
Bảng~\ref{tab:tools} cho thấy: Không ai trong số họ thỏa mãn yêu cầu thời gian thực
trên đây mà không có công việc bổ sung lớn.

\begin{table*}[t]
\centering\small
\caption{So sánh các tùy chọn công cụ so với yêu cầu của kỹ thuật số (không phải
so với một mô phỏng duy nhất).}
\label{tab:tools}
\begin{tabularx}{\textwidth}{@{}l
    >{\centering\arraybackslash}X
    >{\centering\arraybackslash}X
    >{\centering\arraybackslash}X
    >{\centering\arraybackslash}X
    >{\centering\arraybackslash}X@{}}
\toprule
Tiêu chí & SimScale & COMSOL/Maxwell & FEMM & Elmer/FEniCS & \textbf{Python (Được chọn)} \\
\midrule
EM Tần số thấp (Dòng xoáy) & giới hạn & \checkmark & \checkmark & \checkmark & \checkmark (Đường dẫn mất mát có, Lực nâng mở$^\ast$) \\
EM$\to$Liên kết nhiệt & khó & \checkmark & yếu & phức tạp & \checkmark trực tiếp trong bộ nhớ \\
ROM / Khả năng thời gian thực & -- & một phần, đắt tiền & -- & -- & \checkmark \textbf{Thiết kế lõi} \\
Xuất sang Web/AR & -- & -- & -- & -- & \checkmark 1 tệp HTML \\
Hiệu chỉnh cảm biến & -- & độc quyền & -- & tự xây dựng & \checkmark \url{data_io.py} \\
Chạy trên macOS & \checkmark & \checkmark (đắt tiền) & -- chỉ Windows & \checkmark & \checkmark \\
Chi phí & Giới hạn giờ lõi & hàng nghìn €  & miễn phí & miễn phí & miễn phí \\
Hộp đen? & có & có & một phần & không & \textbf{không -- kiểm soát đầy đủ} \\
\bottomrule
\end{tabularx}

\vspace{2pt}
{\footnotesize $^\ast$Đường dẫn mất mát (Dòng xoáy/Mất mát Ohmic) được xác thực;
Lực nâng/Chiều cao lơ lửng vẫn chưa khớp với quan sát -- xem Phần~\ref{sec:limitations}.}
\end{table*}

Hàng cuối cùng của Bảng~\ref{tab:tools} tóm tắt cốt lõi cấu trúc: Các gói FEM
thương mại và nhiều gói mã nguồn mở đóng gói trình giải của họ dưới dạng hộp đen,
bất kể tiêu chí khác -- mỗi thay đổi dòng điện yêu cầu một khởi động lại trình giải
hoàn toàn, thay vì giải một lần được mô tả cộng với tỷ lệ. FEMM được loại bỏ thêm
vì lý do thực tế (chỉ Windows gốc, máy tính phát triển là macOS).

Đối với chạy so sánh độc lập trong một trong các gói thương mại được đặt tên, Hình~\ref{fig:geometry}
và Bảng~\ref{tab:geometry} (Phần~\ref{sec:intro}) có thể được sử dụng trực tiếp: Chúng
chứa cùng các bán kính, số vòng cuộn dây và hằng số vật liệu, trình giải của chính
chúng tôi đọc từ \texttt{params.yaml} -- xây dựng lại là có thể mà không cần liên
hệ với nhóm.


\section{Xác thực và Kết quả Trước đây}
\label{sec:validation}

\begin{table}[t]
\centering\small
\caption{Mất mát sức mạnh của FEM ngoài tuyến ở \emph{biên độ} kích thích $5\,\mathrm A$
và $T_\infty=20\,^\circ$C. Để chuyển đổi thành điểm hoạt động thực tế, xem Ghi chú~\ref{rem:current}.}
\label{tab:losses}
\begin{tabular}{@{}lr@{}}
\toprule
Nguồn mất mát & $P$\,[W] \\
\midrule
Tổn thất dòng xoáy đĩa (Al) & 25{,}81 \\
Tổn thất dòng xoáy lõi sắt + vành & 5{,}86 \\
Tổn thất cuộn dây ohmic (trong 52{,}32 + ngoài 54{,}12) & 106{,}44 \\
\midrule
\textbf{Tổng cộng} & \textbf{138{,}11} \\
\bottomrule
\end{tabular}
\end{table}

Bảng~\ref{tab:losses} cho thấy phân bố mất mát hiện tại. Các bài kiểm tra nội bộ
sau được thực hiện lại mỗi khi thay đổi mô hình; tính toán tính chất của chúng
được cân nhắc một cách cố ý khác nhau:

\begin{itemize}[nosep]
  \item \textbf{Cân bằng năng lượng: $0{,}000\,\%$ độ lệch.} Tại điểm ổn định,
        $\int_V q\,\mathrm dV = \oint_\Gamma h(T-T_\infty)\,\mathrm dA$. Đây là bài
        kiểm tra \emph{có ý nghĩa nhất}, vì nó đặt các số hạng thể tích và biên độc
        lập lại nhau.
  \item \textbf{Kích thước vùng: tất cả độ lệch $<1\,\%$} (lớn nhất:
        $P_{\mathrm{đĩa}}$ với $0{,}767\,\%$) trong so sánh điều kiện biên Dirichlet
        so với Neumann trên vùng $1\times1\,\mathrm m$ -- bằng chứng được cụm khoa
        đặt ra rằng biên ngoài xa đủ.
  \item \textbf{Tỷ lệ $I^2$: $P(2I)/P(I) = 3{,}994$} (giá trị nên $4{,}000$) khi
        giải lại FEM hoàn toàn. Độ lệch còn lại $0{,}15\,\%$ là \emph{không phải}
        lỗi mà là hiệu ứng mong đợi của sự hiệu chỉnh phi tuyến $\mu_r(B)$: Tại
        dòng điện kép, $\mu_r$ trong sắt giảm từ $993$ thành $977$, do đó tổn thất
        tăng lên một chút dưới bình phương.
  \item \textbf{Xác thực chéo: $<10^{-9}$ (tuyệt đối)} giữa bộ tích phân thời gian
        thực được nướng trong JavaScript của chế độ xem AR và bộ tích phân tham chiếu
        Python (\texttt{xval\_twin.py}), được kiểm tra trên bảy hồ sơ dòng điện.
\end{itemize}

\begin{remark}[Kiểm tra $I^2$ trong ROM]
ROM chính nó cung cấp chính xác $4{,}000000$ cho cùng một bài kiểm tra. Tuy nhiên,
giá trị này là \emph{tautological}: ROM \emph{là} theo cấu trúc một phép nhân vô
hướng với $(I/I_{\mathrm{ref}})^2$, do đó không thể thất bại kiểm tra này. Nó chỉ
chứng minh số học chính xác của mức giảm, không phải vật lý. Có ý nghĩa chỉ là giá
trị FEM được đề cập ở trên $3{,}994$.
\end{remark}

\begin{remark}[Quy ước dòng điện]
\label{rem:current}
Giá trị được đo ở giá thử nghiệm $I_{\mathrm{rms}}=5\,\mathrm A$ là giá trị hiệu
dụng; biên độ tương ứng là $\hat I=\sqrt2\,I_{\mathrm{rms}}=7{,}07\,\mathrm A$.
Chuỗi mất mát (Bảng~\ref{tab:losses}) cố ý sử dụng $5\,\mathrm A$ làm biên độ phasor.
Vì tất cả các tổn thất là $\propto I^2$, các số watt được đặt ra ở đó thấp hơn con
số $2$ lần so với sức mạnh thực sự được chuyển đổi ở điểm hoạt động thực tế. Điều
này là cho phép và cố ý, vì hệ số hằng số này được hấp thụ hoàn toàn vào lần phù
hợp của các giá trị dẫn điện lưu thông $hA$ so với nhiệt độ cuộn dây được đo --
dự đoán \emph{nhiệt độ} vẫn còn chính xác. Chuỗi \emph{lực} tương tác với biên độ
thực $\hat I$, vì lực không thể được hiệu chỉnh đi.
\end{remark}

Sự cất giữ bão hòa của sắt cũng không quan trọng, nhưng phải được so sánh chính xác:
Trình giải báo cáo $B_{\max}=0{,}663\,\mathrm T$ như một giá trị \emph{hiệu dụng},
trong khi $B_{\mathrm{sat}}=1{,}5\,\mathrm T$ là một đại lượng vật liệu giá trị đỉnh.
Giá trị đỉnh ngành do đó $\sqrt2\cdot0{,}663 \approx 0{,}94\,\mathrm T$, vì vậy khoảng
$63\,\%$ của $B_{\mathrm{sat}}$. Giả định một $\mu_r$ tuyến tính vẫn có giá trị,
dự trữ tuy nhiên nhỏ hơn, so với một so sánh trực tiếp các giá trị hiệu dụng gợi ý.

Hai điểm hoạt động thực tế được đo bằng máy ảnh IR (Variac-Dial $220^\circ\approx195{,}5\,\mathrm V\to5\,\mathrm A$
làm điểm hoạt động chính, Variac-Dial ở vị trí tối đa $270^\circ$ ($\approx240\,\mathrm V$)
$\to7{,}8\,\mathrm A$ làm điểm lập chỉ số hiệu chuẩn). Mô hình mạng RC của các cuộn
dây được hiệu chỉnh trực tiếp so với bộ tích phân thời gian phi tuyến hoàn toàn và
đạt trạng thái ổn định chính xác ($T_{\mathrm{trong}}=79{,}00\,^\circ$C,
$T_{\mathrm{ngoài}}=74{,}00\,^\circ$C ở $7{,}8\,\mathrm A$). Chỉ các cuộn dây được
sơn đen mới cung cấp các giá trị IR đáng tin cậy (Emissivity không biết của nhôm trần).
Đối với đĩa và lõi, phương pháp được đề xuất không phải là một cảm biến bổ sung mà
là một mảng đo màu đen matte có emissivity già biết ($\varepsilon\approx0{,}95$) ở
dưới đĩa -- không dây, và do đó không phản ứng lại trên lơ lửng.

\textbf{Ổn định xác thực, Transient chưa xác thực.} Sự phân biệt này là quan trọng:
Trạng thái ổn định chính xác theo cấu trúc, vì $hA_{\mathrm{trong/ngoài}}$ chính xác
được điều chỉnh với nó. Quá trình thời gian tuy nhiên \emph{không} được xác thực
với nó. Chống lại phép đo transient duy nhất có sẵn -- một gradual dòng điện hai bước
($5\,\mathrm A$ cho đến $t=300\,\mathrm s$, sau đó $7{,}75\,\mathrm A$) với năm
dấu thời gian của nhiệt độ cuộn dây ngoài -- mô hình hiện tại cho kết quả lỗi
$\mathrm{RMS}=11{,}5\,^\circ$C (Bảng~\ref{tab:transient}): Nó ấm lên một cách có
hệ thống quá chậm. Nguyên nhân là một thay đổi mô hình được chấp nhận cố ý; nó
được thảo luận trong Phần~\ref{sec:limitations}.

\begin{table}[t]
\centering\small
\caption{So sánh transient chống lại gradual dòng điện được đo (Cuộn ngoài,
$T_\infty=29\,^\circ$C). Trạng thái ổn định khớp chính xác, quá trình ấm lên chưa.}
\label{tab:transient}
\begin{tabular}{@{}rrrr@{}}
\toprule
$t$\,[s] & Được đo\,[$^\circ$C] & Mô hình\,[$^\circ$C] & $\Delta$\,[K] \\
\midrule
140 & 40{,}00 & 33{,}55 & $-6{,}45$ \\
200 & 43{,}00 & 34{,}70 & $-8{,}30$ \\
300 & 48{,}50 & 36{,}22 & $-12{,}28$ \\
360 & 55{,}75 & 40{,}00 & $-15{,}75$ \\
450 & 56{,}00 & 43{,}55 & $-12{,}45$ \\
\midrule
\multicolumn{3}{@{}l}{\textbf{RMS}} & \textbf{11{,}53} \\
\bottomrule
\end{tabular}
\end{table}

Để phân bổ vai trò giữa phép đo và mô hình, nó được giữ rõ ràng: Giá thử nghiệm
có \emph{không} cảm biến nhiệt độ được lắp đặt vĩnh viễn. Được đo chỉ là dòng điện
cuộn dây; nhiệt độ là đại lượng mà mô hình \emph{tính toán} từ nó (Phần~\ref{sec:realtime_input}).
Giá trị đo nhiệt độ để đối chiếu do đó đến từ ảnh IR thủ công, không phải từ nhật ký
liên tục. Hiệu chuẩn lại tự động của mức độ truyền nhiệt từ nhật ký cảm biến do đó
không thể được thực hiện trên thiết lập này -- một giới hạn được chấp nhận cố ý, không
phải một khoảng trống triển khai mở (xem Phụ lục~\ref{app:sensor}).


\section{Kết quả Dự kiến và Triển vọng}
\label{sec:outlook}

Dự kiến gần thời gian: bắt đầu hoạt động phần cứng cảm biến và do đó chuyển từ
giả định không đổi $I=5\,\mathrm A$ sang một \textbf{dòng điện được đo thực sự $I(t)$},
nuôi dưỡng kỹ thuật số nhiệt lúc thực thi (Phần~\ref{sec:realtime_input}); từ đó
một \emph{và} đường cong làm lạnh ấm lên thực tế, khiến mạng RC chưa xác định
(đặc biệt là nút không khí chung và phân chia hai nút của cuộn dây) xác định lần
đầu tiên và do đó sẽ đóng độ lệch transient hiện đang mở (Phần~\ref{sec:limitations}) --
phía nhiệt độ của các đường cong này được ghi lại không tiếp xúc bằng máy ảnh IR,
với mảng đo màu đen matte như một điều trị cho sự không chắc chắn emissivity; cũng
như một chế độ xem AR được xuất bản với Mã QR, có thể gọi mà không cài đặt. Trung hạn:
bảng hiệu chuẩn denser phím quay Variac $\to$ dòng điện và xác định quyết đoán hơn của
độ thẩm thấu sắt để làm rõ câu hỏi lực nâng được đặt tên trong Phần~\ref{sec:limitations}.


\section{Câu hỏi Mở và Giới hạn}
\label{sec:limitations}

Trong ý nghĩa của sự minh bạch khoa học: Chiều cao lơ lửng dự đoán ($z_{\mathrm{eq}}\approx11{,}7$--$14{,}7\,\mathrm{mm}$)
không khớp với khoảng cách quan sát được ($7$--$8\,\mathrm{mm}$); bão hòa từ tính
được loại bỏ làm nguyên nhân (Thay đổi lực $<0{,}1\,\%$), nguyên nhân có khả năng
nhất là giá trị giữ chỗ $\mu_r$ không được đo của sắt. Điều chỉnh ngược tùy ý của
$\mu_r$ để khớp với chiều cao quan sát được cố ý \emph{không} được thực hiện. Chuẩn
TEAM-28 gốc (không sắt, $20\,\mathrm A$) cũng sai lệch $\approx28\,\%$ từ bảng tham
chiếu -- tạm dừng, không chặn.

\textbf{Độ lệch Transient của cuộn dây.} Lỗi $\mathrm{RMS}=11{,}5\,^\circ$C được
ghi lại trong Bảng~\ref{tab:transient} là độ lệch mở lớn nhất hiện tại của mô hình
nhiệt. Nguyên nhân của nó được biết đến và có thể hiểu được: Cuộn dây được mô hình
hóa như một cơ thể hai nút (bề mặt nhìn thấy IR cộng với phần trong cuộn trễ).
Trong phiên bản mô hình trước đó, toàn bộ sức mạnh mất mát ohmic được đặt trên nút
bề mặt, chỉ mang $22{,}4\,\%$ dung lượng nhiệt -- điều đó có kết quả là một tỷ lệ
ấm lên tốt để đo lường, nhưng về mặt vật lý không có cơ sở. Sức mạnh mất mát bây
giờ được phân phối trong tỷ lệ khối lượng trên cả hai nút, để lại trạng thái ổn định
chứng minh không thay đổi, nhưng tỷ lệ ấm lên ban đầu bằng một yếu tố $\approx4{,}5$.
Mô hình kể từ đó là nhất quán hơn về mặt vật lý, nhưng phù hợp với gradual được đo
tệ hơn. Các tham số ngành (phân chia dung lượng \texttt{coil\_C\_scale}, ghép nối
nội bộ \texttt{coil\_G\_wind}) không thể xác định duy nhất từ phép đo ấm lên thuần
-- chúng yêu cầu thêm một \emph{đường cong làm lạnh}, vẫn đang chờ. Cho đến lúc đó
áp dụng: trạng thái ổn định được xác thực, quá trình thời gian rõ ràng không.

Không có đường cong làm lạnh này, nút không khí chung của mạng RC vẫn còn xác định
dưới. Cũng mở: Trong mô hình đĩa và cuộn dây đạt gần như cùng một nhiệt độ, trong
khi dữ liệu IR gợi ý một cuộn dây rõ ràng nóng hơn; các quyết định chỉ có thể được
thực hiện với một bài kiểm tra đáng tin cậy ở đoạn dưới của đĩa, tức là với mảng
đo màu đen matte được đề cập ở trên. Cuối cùng, phản hồi $\sigma(T)$ của cuộn dây
hiện tại được chứa ngầm chỉ qua lần phù hợp $hA$ (Ghi chú~\ref{rem:sigma}).


\section{Tóm tắt}
\label{sec:conclusion}

Dự án theo một con đường có cơ sở vật lý cho một kỹ thuật số thời gian thực: Một
lần giải FEM độ phân giải cao được giảm xuống thành thời gian chạy mili giây thông
qua tỷ lệ $I^2$ và một mô hình giảm bậc có thể hiệu chỉnh được -- một cấu trúc không
thể thực hiện được với công cụ CAE thương mại được đóng gói. Mỗi thay đổi mô hình
được phát hành độc quyền thông qua tiêu chí vật lý tự động cứng (cân bằng năng lượng,
Nhất quán $I^2$, xác thực chéo, điều chỉnh dữ liệu cảm biến thực). Các kết quả hiệu
chuẩn đầu tiên trên nhiệt độ cuộn dây xác nhận khả năng của mô hình; chiều cao lơ
lửng vẫn còn là một câu hỏi vật lý mở và sẽ được điều tra thêm với phần cứng cảm biến
được lên kế hoạch.


%%%%%%%%%%%%%%%%%%%%%%%%%%%%%%%%%%%%%%%%%%%%%%%%%%%%%%%%%%%%%%%%%%%%%%%%%%%%%%%%
\appendices

\section{Bản sao Đầy đủ của các Phương trình Vật lý}
\label{app:derivation}

Dạng yếu đối xứng trục của phương trình dẫn nhiệt với trọng lượng thể tích
$2\pi r$ đọc
\begin{multline*}
  \int k\,\nabla T\!\cdot\!\nabla v \cdot 2\pi r \,\mathrm dA
  + \oint_\Gamma h\,T\,v\cdot 2\pi r\,\mathrm ds \\
  = \int p\,v\cdot 2\pi r\,\mathrm dA
  + \oint_\Gamma h\,T_\infty\,v\cdot 2\pi r\,\mathrm ds .
\end{multline*}
Đối với các phần tử P1, $\nabla T$ là phần tử hằng số, do đó ma trận độ cứng
phần tử $K_e = 2\pi k\,(bb^\top+cc^\top)\,A\,r_c$ tuân theo trực tiếp từ các
độ dốc của hàm hình $b,c$, diện tích phần tử $A$ và bán kính trung tâm $r_c$.
Ma trận biên Convection với cân nặng bán kính tuyến tính kết quả
$K_{\mathrm{edge}} = 2\pi h\,\tfrac{L}{12}
\left(\begin{smallmatrix} 3r_a+r_b & r_a+r_b \\ r_a+r_b & r_a+3r_b\end{smallmatrix}\right)$.
Trong \emph{vấn đề nhiệt} con, trục $r=0$ là một điều kiện đối xứng và vẫn là
điều kiện biên tự nhiên chưa được áp dụng (các cạnh biên trên trục được bỏ qua
trong lắp ráp Convection). Trong \emph{vấn đề điện từ} con, điều đó \textbf{không}
được áp dụng: Ở đó $A_\varphi=0$ được đặt rõ ràng trên trục là điều kiện Dirichlet,
vì thành phần phương vị của thế năng vectơ phải biến mất về mặt vật lý như $r\to0$.
Cho phần tùy này kết quả tương tự một ma trận độ cứng từ Phương trình~\eqref{eq:em}.
Số hạng $j\omega\sigma$ làm cho nó phức tạp; nó được lắp ráp như một ma trận
phức tạp duy nhất và được phân tích trực tiếp trong số học phức tạp (không tách
thành hệ thống bộ phận thực và tưởng tượng). Mật độ từ thông từ tính kết quả
$B_r=-\partial A_\varphi/\partial z$, $B_z=\tfrac1r\,\partial(rA_\varphi)/\partial r$.

Về mặt vật lý, Phương trình~\eqref{eq:em} mô tả trường AC: Số hạng $j\omega\sigma
A_\varphi$ là phản ứng dòng xoáy trong mỗi cơ thể dẫn điện (đĩa, sắt); số hạng
$\nu A_\varphi/r^2$ là hiệu ứng hình học của đối xứng trục và lý do tại sao 2D
thay vì 3D là đủ. Ở $I_{\mathrm{rms}}=5\,\mathrm A$ qua cuộn quay $1000$-cuộn
ngoài, $B_{\max}$ trong lõi sắt $0{,}663\,\mathrm T$ như một giá trị hiệu dụng,
tương ứng một giá trị đỉnh $\approx0{,}94\,\mathrm T$ -- so với $B_{\mathrm{sat}}=1{,}5\,\mathrm T$
do đó chưa bão hòa (so sánh Phần~\ref{sec:validation}).

\section{Đường ống Phần mềm Đầy đủ}
\label{app:software}

Bảng~\ref{tab:modules} trả lời rõ ràng thư viện nào giải phương trình nào bằng
phương pháp số nào -- bổ sung để Hình~\ref{fig:pipeline} (Luồng dữ liệu) và
Phụ lục~\ref{app:derivation} (Bản sao của các phương trình tự chúng).

\begin{table}[h]
\centering\scriptsize
\caption{Mô-đun cốt lõi của đường ống phần mềm: phương trình được giải và phương
pháp số cho mỗi tệp.}
\label{tab:modules}
\begin{tabularx}{\linewidth}{@{}p{1.6cm} X X@{}}
\toprule
Tệp & Giải & Phương pháp \\
\midrule
\url{config.py} & Tham số tải/bình thường hóa (mm$\to$m), $I_{\mathrm{peak}}=I_{\mathrm{rms}}\sqrt2$ & PyYAML; Lớp dữ liệu, không giải phương trình \\
\url{em_solver.py} & Eq.~\eqref{eq:em} $\to$ Eq.~\eqref{eq:loss}, Nâng lực & \url{scipy.sparse.linalg.spsolve} (phức tạp, LU) \\
\url{thermal_solver.py} & Eq.~\eqref{eq:fourier}, ổn định & \url{scipy.sparse.linalg.spsolve} (thực, SPD) \\
\url{rom.py} & $\tau{=}C/UA$ từ $\hat P,\Delta\hat T$; ODE xác thực & \url{scipy.integrate.solve_ivp} \\
\url{twin_core.py} & Mạng RC (Cuộn dây/Sắt/Không khí); Lơ lửng & Euler rõ ràng, $n_{\mathrm{sub}}{=}\lceil\Delta t/(0{,}05\tau)\rceil$; Lơ lửng giải quyết gần \\
\url{refit_hA.py} & Không của $T_{\mathrm{ss}}(hA)-T_{\mathrm{meas}}$ & \url{scipy.optimize.fsolve} \\
\url{xval_twin.py} & -- (Xác minh phần mềm) & Playwright: Integrator Python- vs.\ JS, $|\Delta|{<}10^{-9}$ \\
\url{build_twin_html_fem.py} & nướng: Công thức giống hệt nhau trong JavaScript & -- \\
\url{data_io.py} & Cầu cảm biến $\to$ \url{TwinState.step(I,dt)} & Phân tích chuỗi/Mock-CSV (NumPy) \\
\bottomrule
\end{tabularx}
\end{table}

Theo thuật toán, đường ống chia chính xác thành hai bước từ Phần~\ref{sec:concept}:
\textbf{Ngoại tuyến} \url{em_solver.py} và \url{thermal_solver.py} mỗi cái giải một
hệ phương trình tuyến tính \emph{một lần trực tiếp} (Phân rã LU, không lặp lại) và
cung cấp $\hat q(r,z)$, $\hat P_{\mathrm{Cu}}$ cũng như $\tau$. \textbf{Trực tuyến}
\url{twin_core.py} tích hợp mạng RC kết quả với một thủ tục Euler rõ ràng được viết
tay -- được chọn vì tự do thư viện (\texttt{numpy+stdlib only}) duy nhất cho phép
xuất sang JavaScript sau; kích thước bước được chia tự động trên mỗi cuộc gọi thành
$n_{\mathrm{sub}}$ bước con không quá $5\,\%$ của $\tau$, để giữ cho quy trình được
ổn định về mặt số học ngay cả khi $\Delta t$ lớn (ví dụ: thời gian nhanh). Ngoại lệ
duy nhất là động lực lơ lửng: phương trình chuyển động (Oscillator lò xo-khối lượng
được giảm chấn) của nó có một giải pháp phân tích chính xác và do đó được đánh giá
trực tiếp trên mỗi cuộc gọi, không được tích hợp lặp lại.

\section{Thiết lập Đo lường cho Đầu vào Dòng điện}
\label{app:sensor}

Đặc tả của khoa là thiết kế một giải pháp đo lường nhỏ dựa trên Arduino. Được triển
khai là cố ý \emph{chính xác một} đại lượng được đo: dòng điện cuộn dây. Thiết lập
bao gồm một cảm biến dòng điện Hall ACS712-20A trên Arduino UNO (R3 hoặc R4;
pincompatible, firmware giống nhau), gửi các dòng CSV \url{millis,I_rms_A} qua USB
mỗi giây.

\textbf{Tại sao chỉ dòng điện?} Vì dòng điện là \emph{đầu vào} của kỹ thuật số và
nhiệt độ là \emph{đầu ra} của nó. Giá thử nghiệm không có cảm biến nhiệt độ được lắp
đặt, và một nhiệt điện được gắn vào đĩa lơ lửng sẽ làm phiền một cách chính xác kích
thước được đo lường. Phản ứng nhiệt độ do đó không tiếp xúc qua máy ảnh IR. Điểm yếu
đã biết của nó -- nhôm trần có emissivity thấp không biết và cung cấp các giá trị
không sử dụng được (xác nhận 2026-06-23) -- không được giải quyết bằng cảm biến liên
hệ, mà là một mảng đo màu đen matte có emissivity đã biết trên phía dưới của đĩa.

Danh sách hóa chất: Mô-đun ACS712-20A; Tụ điện gốm $0{,}1\,\mu$F (OUT$\to$GND, bộ
lọc nhiễu); Các clamp vít hai cực cho phía điện áp mạng; Jumpers. \textbf{An toàn:}
IP+/IP$-$ nằm trên mạch điện áp mạng $0$--$240\,\mathrm V$ và được cách ly từ
phần tín hiệu $0$--$5\,\mathrm V$; cả hai phía không được bao giờ được cầu nối, và
Variac phải được đặt trên $0$ trước bất kỳ thay đổi dây nào.

Bắt đầu vận hành: (1) Kết nối Arduino mà không kết nối mạng -- dòng chẩn đoán phải
báo cáo một sự sai lệch DC gần $2{,}5\,\mathrm V$ và một dòng điện gần $0\,\mathrm A$
(cơ sở nhiễu); (2) Variac ở $0$, các hành động tại IP+/IP$-$ clamp; (3) Cách từ từ
đến vị trí Dial $220^\circ$ (điểm hoạt động chính, $\approx5\,\mathrm A$) và kiểm tra
giá trị dòng điện được báo cáo so với đồng hồ vôn kế, từ đó xác định hệ số $k$ từ
Phương trình~\eqref{eq:irms}; (4) Ghi lại một lần chạy trong $\geq3\tau\approx20\,\mathrm{min}$
và chạy kỹ thuật số với nó; (5) Đồng thời ghi lại ảnh IR ở các khoảng thời gian cố
định như một ngành cơ sở độc lập của nhiệt độ được tính toán -- bao gồm một pha làm
lạnh sau khi tắt, yêu cầu cho câu hỏi transient được đặt tên trong Phần~\ref{sec:limitations}.

\bibliographystyle{IEEEtran}
\bibliography{dt4tm_references}

\end{document}
